\section[Tests first]{Tests first}
\sectionframe{06}{Tests first}
  {How does a test written before the code keep an agent on track?}
  {\faCheckDouble}

% ---------------------------------------------------------------- the loop
\begin{frame}[source={Beck (2002), Test-Driven Development: By Example}]{Test-Driven Development (TDD)}
  \begin{tikzpicture}[x=1mm,y=1mm,
      phase/.style={box,text width=28mm,minimum height=13mm,font=\footnotesize,line width=0.8pt},
      gain/.style={inner sep=0pt,font=\footnotesize,text=udsdark,align=left},
      tick/.style={inner sep=0pt,font=\footnotesize,anchor=base east}]
    \useasboundingbox (0,-1.36) rectangle (144,-55);
    \node[phase,draw=udsred,fill=redtint] (red) at (72,-8)
      {{\color{udsred}\bfseries\small RED}\\[0.4mm] write a failing test for one behaviour};
    \node[phase,draw=testgreen,fill=greentint] (green) at (96,-40)
      {{\color{testgreen}\bfseries\small GREEN}\\[0.4mm] write just enough code to pass it};
    \node[phase,draw=udsblue,fill=bluetint] (ref) at (48,-40)
      {{\color{udsblue}\bfseries\small REFACTOR}\\[0.4mm] tidy the code; every test stays green};
    \draw[flow] (red.east) to[bend left=35] (green.north);
    \draw[flow] (green.west) -- (ref.east);
    \draw[flow] (ref.north) to[bend left=35] (red.west);
    \node[lbl,align=center] at (72,-27) {one behaviour\\per cycle};
    % what each step gives us; each tick sits on its gain's first baseline (with a text width,
    % a node's base is its first line); the two gains of RED start at x = 98, clear of the arrow
    % from RED to GREEN
    \node[gain,anchor=west] (g1) at (98,-5.2) {the requirement, written first};
    \node[gain,anchor=west] (g2) at (98,-10.8) {proof that the test can fail};
    \node[gain,anchor=west,text width=20mm] (g3) at (118.4,-40) {small steps,\\easy to review};
    \node[gain,anchor=west,text width=22mm] (g4) at (8,-40) {a regression net\\for later changes};
    \foreach \g/\c in {g1/udsred,g2/udsred,g3/testgreen,g4/udsblue}
      \node[tick,text=\c] at ([xshift=-1.4mm]\g.base west) {\faCheck};
    \node[gain,anchor=base] (g5) at (74,-53.4) {tests call clear interfaces: the seams};
    \node[tick,text=udsdark] at ([xshift=-1.4mm]g5.base west) {\faCheck};
  \end{tikzpicture}
  \note{\textbf{25:00--26:00}\par
    Test-driven development in Kent Beck's formulation is a short loop.\par
    Red: write a test for one behaviour and watch it fail.\par
    Green: write just enough code to make it pass.\par
    Refactor: tidy the code while all tests stay green.\par
    Then the next behaviour.\par
    What we get: the requirement written down as a test before the code exists, proof that each
      test can fail, small steps, a regression net, and pressure towards clear interfaces, the
      seams where tests observe behaviour.}
\end{frame}

% ------------------------------------------------------- why for an agent
\begin{frame}[source={Anthropic (2026), Claude Code best practices · Mathews \& Nagappan (2024), ASE ’24}]{Why tests first help an agent}
  \begin{tikzpicture}[x=1mm,y=1mm,
      node2/.style={box,minimum height=9mm,inner sep=1.6mm},
      item/.style={anchor=base west,inner sep=0pt,font=\footnotesize,text=udsdark},
      head/.style={anchor=base west,inner sep=0pt,font=\footnotesize\bfseries}]
    \useasboundingbox (0,0) rectangle (144,-54);
    % the agent edits the code and runs the tests
    \node[node2,artifact,font=\ttfamily\footnotesize\bfseries,minimum width=20mm] (src) at (12,-8) {src/};
    \node[node2,agent,font=\footnotesize\bfseries,minimum width=26mm] (ag) at (62,-8) {\faRobot\enspace agent};
    \node[node2,artifact,draw=testgreen,fill=greentint,font=\ttfamily\footnotesize\bfseries,minimum width=20mm]
      (tst) at (114,-8) {tests/};
    \draw[flow] (ag.west) -- node[lbl,above] {edits} (src.east);
    \draw[flow] (ag.east) -- node[lbl,above] {runs} (tst.west);
    \node[anchor=south,inner sep=0pt,font=\scriptsize,text=testgreen] at ([yshift=1mm]tst.north)
      {\faLock\enspace approved and committed first};
    % pass or fail comes back: the signal
    \draw[flow,draw=testgreen,line width=0.8pt] (tst.south) to[out=-90,in=-90,looseness=0.55]
      node[pos=0.5,below,inner sep=0.8mm,font=\scriptsize\bfseries,text=testgreen] {pass or fail} (ag.south);
    % what the signal gives the agent
    \node[head,text=testgreen] at (0,-28.4) {what it gives the agent};
    \foreach \t [count=\i] in {a check it runs itself,{the intent, fixed first},one small goal per cycle,
        earlier decisions still enforced}
      \node[item] at (0,-28.4-\i*4.7) {\textcolor{testgreen}{\faCheck}\enspace\t};
    % how it gets gamed: aimed at the tests, and what catches it
    \node[head,text=udsred] (gamed) at (70,-28.4) {how it gets gamed};
    \foreach \t [count=\i] in {tests edited or skipped,expected values copied from the code,
        a mock that replaces the logic,tests fitted after the code}
      \node[item] at (70,-28.4-\i*4.7) {\textcolor{udsred}{\faTimes}\enspace\t};
    \draw[udsred,line width=0.6pt,dashed,-{Stealth[length=1.7mm,width=1.3mm]}]
      ([xshift=1.5mm]gamed.east) to[out=0,in=-90] (121,-12.6);
    \node[item,text=udsblue] at (70,-53)
      {caught by: test diffs reviewed first · you rerun every test};
  \end{tikzpicture}
  \note{\textbf{26:00--27:00}\par
    Anthropic's guidance for Claude Code puts it plainly: without a check it can run, an agent
      stops when the work looks done.\par
    A failing test turns done into a signal the agent reads and iterates on.\par
    Written first, from the spec, the test records intent instead of mirroring the
      implementation, and the suite keeps earlier decisions enforced in later sessions.\par
    The failure modes are real: agents weaken, skip or edit tests, recompute expected values the
      way the code does, or replace the logic under test with a mock, a stand-in that returns
      whatever the test expects.\par
    Hence our guardrails: read the changes to the tests first, and rerun every test yourself
      rather than trust the agent's report.\par
    A benchmark study found that giving the model the tests with the task solved more problems;
      that is evidence for code generation, not for research code.}
\end{frame}
